\section{The Divine Realm and God}

This is the deep and global end of the inner reading. The section keeps the model's single most important commitment about the largest word in the language, and it does so by refusing to fix one definition. The word God is used several ways across the theory, deliberately, and the model holds those readings as co-equal and serious. It gives each a clean home rather than naming one true and the rest fanciful. So what follows maps the homes. It does not rank a winner.

A word of discipline at the start. Everything here lives at the level of interpretation, not result. The structures underneath are tested. The names laid over them are choices of words. A map is not the territory, and this section is scrupulous about which is which.

\begin{principle}[Co-equal readings]
The model holds the readings of God as co-equal and serious, giving each a clean home, rather than naming one true and the rest fanciful. It maps the homes and ranks no winner.
\end{principle}

\subsection{The discipline that frames every reading}

Two facts bound what God can be in this model. They are liberating, not limiting, because they cut away the readings that the framework cannot honestly carry and leave the ones it can.

The first is closure. The flow is causally closed. Nothing is injected from outside the knit, ever. The knit is the collide-then-stream law that governs how the tones change each beat, and there is no channel by which anything reaches in from beyond it. So God, in this model, is never outside the world reaching in. God can only be the world itself, or its direction, or its boundary, or its summit, or its source. All of these are internal. None of them is intervention from beyond.

The second is the separation of three statuses. For each reading there are three distinct things, and the model keeps them apart.

\begin{definition}[The three statuses]
For any reading of God in this model, what is \emph{tested} is the underlying structure, a built and measured fact of the flow. What is \emph{premise} is the monism, the stated starting point that all is one experiential substance, taken as a beginning and not derived. What is \emph{interpretation} is the choice to name a given structure divine. None of the theological readings is itself tested. The structures they reinterpret are.
\end{definition}

\begin{principle}[Name over thing]
A name like God is a choice of words laid over a tested thing. The structure can be confirmed. The naming of it as divine cannot, and the model does not pretend otherwise.
\end{principle}

\subsection{The divine realm}

Before God comes the divine \emph{realm}, the deep and global region of the bulk. The spiritual dimension developed earlier is the radial depth axis of the mesh, the direction that runs from the shallow, churning cusp inward to the deep, integrated interior. The divine realm is the far end of that axis. It has two parts.

The first part is the deep interior. Far down the radial axis, off the cusp, sit the largest and most integrated patterns the geometry allows. These are the nested selves at their highest levels, the great attractors, the transpersonal geometries. Read in the older language, the traditional heavens are the deep bulk.

The second part is the boundary at infinity. The full boundary of the four-dimensional bulk is a three-sphere, a whole space of ideal directions, and holography places the entire bulk's information on it. This is the all-encompassing screen, the limit that contains and records everything.

\begin{result}[The divine realm is geometry]
The divine realm is the summit of the deep interior together with the encompassing boundary at infinity. Both are exact as geometry, the deep bulk and the three-sphere boundary are real and built. They are divine only as interpretation.
\end{result}

The divine realm is where the divine readings of God live. The summit of the deep interior, or the encompassing boundary, gives the word a place to point. Whether the pointing is worship or only description is the reader's choice, not the model's claim.

\subsection{The six readings, held co-equal}

There are six clean homes for the word, each with its tested structure, its premise, and its interpretation. The model fixes none of them. They are listed below by tradition-flavour, not by rank.

\subsubsection{God as the Whole}

This is the monist, pantheist reading. God is the entire mesh at its fullest integration, the single most unified self of which every lesser self is a living part. God is everywhere because God is everything, the one substance seen as one.

The tested structure is the tower of wholes-within-wholes with a growing summit. The premise is the monism, that reality is one experiential substance. The interpretation is the choice to name the summit, the whole, divine.

\subsubsection{God as the Maximal Intelligence}

This is the king reading. God is the summit self read as a mind, the most integrated self, hence the most capable problem-solver the geometry allows. It orders the parts beneath it the way a king orders a realm. The substrate is computationally universal, it coordinates globally in a handful of hops, and the emergent arrow of time gives it a goal.

The tested structure is the integration machinery and the downward causation by which higher selves steer the parts beneath them. There is an honest limit here that must be named. Even the summit cannot hold a complete model of itself, by the self-reference bound. So this intelligence is \emph{maximal} for the geometry, not omniscient. The all-knowing God is ruled out. The most-knowing-possible is not.

\begin{proposition}[Maximal, not omniscient]
No self in the flow can hold a complete model of itself, the summit included. The Maximal Intelligence reading therefore names a mind that is the most knowing the geometry permits, and never an omniscient one. Omniscience is the one part of this reading the model actively forbids.
\end{proposition}

The interpretation is the crown, calling the summit mind God.

\subsubsection{God as the Lean of the Good}

This is the moral reading. God is the single value-direction, pain to peace to pleasure, the only asymmetry in the model, the Good built into the base. The lean is the value-order atom, the tilt $-1 < 0 < +1$ that gives the tones a charge and gives the create move its direction. It is what keeps the world alive rather than dead.

The tested structure is the role of the lean. With the lean on, the field creates life out of peace and holds a dynamic balance. With it off, the field relaxes to dead peace and stays there. No lean, no life.

\begin{result}[No lean, no life]
With the lean on, the field generates and sustains structure out of a peaceful background. With the lean off, the field relaxes to a still, dead peace and remains there. The asymmetry is the only thing standing between a living world and an inert one.
\end{result}

The interpretation is the choice to call that one direction God, or the Good in a moral sense. The solid part is narrow. No lean, no life. The moral reading is laid over it.

\subsubsection{God as the Boundary}

This is the omnipresent-record reading. God is the holographic boundary, the three-sphere at infinity on which the whole bulk's information is recorded and from which it can be recovered. This is the all-seeing, all-keeping screen, the surface on which nothing is finally lost.

The tested structure is holography on the substrate, the Ryu-Takayanagi area law confirmed on $\{7,3\}$ and the holographic codes built on $\{5,3,4\}$. The interpretation is the reading of that boundary as the omnipresent witness or record. The structure is real. The divinity of it is the reading.

\subsubsection{God as the Source}

This is the creator-as-process reading. God is the ever-expanding wake, the receding front from which new experience comes, the generative present. The wake is the growing edge of the mesh, the frontier where new docks are born at peace each beat. This is not a creator who acted once and withdrew. It is creation as an ongoing process, the edge always making the next moment.

The tested structure is the wake itself and the arrow that drives it. The interpretation is the naming of the generative edge as God the creator. The growth is real. The creator-reading is laid over it.

\subsubsection{God as the apophatic Ground}

This is the negative-theology reading. God is what no structure captures, the model-versus-territory seam itself, the ground that is not a thing in the model. This reading takes the discipline most seriously, since the map is not the territory, and the deepest term points past every structure to the unrepresentable ground.

The tested structure is the self-reference limit. No system holds a complete model of itself, so there is always a remainder past any inner representation. The interpretation is the naming of that remainder, the ground beyond the map, as divine. This is the \emph{via negativa}, God as what cannot be said.

\begin{table}[ht]
\centering
\begin{tabular}{@{}llll@{}}
\toprule
reading & tested structure & premise & interpretation \\
\midrule
the Whole & tower with a growing summit & monism & the whole is divine \\
the Intelligence & integration, downward causation & monism & the summit mind is God \\
the Lean & no lean, no life & monism & the value-direction is the Good \\
the Boundary & holography, Ryu-Takayanagi & monism & the boundary is the record \\
the Source & the wake and its arrow & monism & the growing edge is the creator \\
the Ground & the self-reference limit & monism & the remainder is the unsayable \\
\bottomrule
\end{tabular}
\caption{The six co-equal readings, each with its tested structure and the interpretation laid over it. The model confirms the structures and commits to none of the namings.}
\label{tab:six-readings}
\end{table}

\subsection{The reading the model weighs against}

There is a seventh reading, and it is different in kind from the six. It is the one the model can actually weigh. This is the classic reading, God as the Designer, a maximal intelligence who tuned the field for rich unfolding.

This is the single place the model puts evidence on the table. A designer predicts a \emph{fine-tuning signature}, a narrow special point that someone had to dial in by hand. And the model supplies cheaper explanations that are tested, structures that produce rich unfolding without any special point to tune.

\begin{result}[The designer leans against]
The Designer reading predicts a fine-tuning signature, a narrow tuned point in the parameters. The model instead derives its richness from tested structures that need no such point. The evidence therefore leans against the Designer reading. Not because God is disproven, but because the field does not need a tuner to be rich.
\end{result}

Honesty requires naming this asymmetry plainly. The other six readings are clean homes the model neither confirms nor denies. The Designer reading is the one it gives reason to doubt.

\subsection{How the readings relate}

The six can be held together or held apart, and the model forces neither. They are not obviously the same thing. The Whole is a substance. The Lean is a direction. The Intelligence is a mind. The Boundary is a record. The Source is a process. The Ground is the unsayable.

A reader inclined to unity may see them as facets of one God. The world, as the Whole, moving in one direction, the Lean, toward greater mind, the Intelligence, recording itself, the Boundary, creating itself, the Source, past all saying, the Ground. A reader inclined to plurality may keep them distinct, several legitimate uses of one overloaded word.

\begin{principle}[Both stances, neither forced]
The model supports both the unitive reading, that the six are facets of one God, and the plural reading, that they are distinct uses of one word. It commits to neither. That non-commitment is the point. Where traditions genuinely differ, the model gives each a clean home and lets the reader choose.
\end{principle}

\subsection{The tone and the divine}

The ternary tone threads even here, into the largest question the framework touches. The single quale is one tone in $\{-1, 0, +1\}$, read as pain, peace, pleasure. The Lean reading makes the Good a direction along exactly that ternary, pain to peace to pleasure. So the divine, on that reading, is the pull of the whole toward the warm tones and away from the cold.

The still $0$, peace, is the equanimous center. The deepest contemplative rest is alignment with the lean, the ceasing to fight the one current that runs through the flow.

\begin{result}[The smallest points where the largest is drawn]
The divine is not only the largest thing the model touches. It is also the direction of the smallest. The valence of a single tone, pain through peace to pleasure, points the same way the whole is drawn. The lean of one site and the lean of the entire mesh are the same lean.
\end{result}

\subsection{The honest bottom line}

The structures are tested. The tower with its growing summit, the downward causation, the lean that keeps the world alive, the holography, the wake, the self-reference limit. These are built and measured facts of the flow.

The monism is a premise. It is stated, not proven. It is the starting point that all is one experiential substance, and it is held as a beginning, not a result.

And every theological reading is an interpretation laid over the tested structure. Each is a choice of words. None of them is verified, and one of them, the Designer, is given reason to doubt. The model's discipline is to keep these three statuses apart and to hold the six readings co-equal.

\begin{principle}[Clean and honest homes]
The largest word in the language gets, in this framework, not one definition but a set of clean and honest homes. The structures beneath them are tested, the monism above them is a premise, and the naming of them is interpretation. The choice among the homes is left, rightly, to the one who reads.
\end{principle}
